\chapter{Status and Open Issues}
\label{ch:status}

This chapter lists what the reference simulator implements, where it disagrees with the older design notes on the
project wiki or with itself, and the design questions that remain. It is the list to read before relying on anything
else in the manual.

\section{Implementation Status}

\begin{table}[H]
	\centering
	\begin{tabular}{p{5.8cm}p{3.4cm}p{5.2cm}}
		\toprule
		\textbf{Feature}                                   & \textbf{Status}                                              & \textbf{Notes}                                                                  \\
		\midrule
		Registers, palette and VRAM access                 & Implemented                                                  & Open bus reads zero                                                             \\
		Package, pins, signal timing, video output         & Specified                                                    & The simulator has no pins; Chapters~\ref{ch:package} and~\ref{ch:signal-timing} \\
		Mode 1 backgrounds, 8 by 8 tiles, 32 by 32 maps    & Implemented                                                  & Layer order corrected with the sprite prototype                                 \\
		Mode 1 scrolling (specified)                       & Not yet implemented                                          & Registers stored                                                                \\
		Mode 1 16 by 16 tiles, double tilemaps (specified) & Not yet implemented                                          & Registers stored                                                                \\
		Orthogonal tile flips (specified), priority bit    & Not yet implemented; priority in Mode 3 and the status layer &                                                                                 \\
		2-bit tile data (Mode 1)                           & Not yet implemented                                          &                                                                                 \\
		Graphics disable, brightness (specified)           & Not yet implemented                                          & Bits stored                                                                     \\
		Background disables                                & Modes 2 and 3 only                                           &                                                                                 \\
		Sprites, sprite disables and sizes                 & Prototype                                                    & Proposed, both modes; Chapter~\ref{ch:sprites}                                  \\
		Vertical sync interrupt (NMI)                      & Implemented                                                  & At line 224; on level 4, not the NMI yet                                        \\
		Vertical blanking flag and current line register   & Implemented                                                  & \texttt{VBL} in Status, \texttt{VLINE} at \texttt{0x0018}                       \\
		PAL timing and the output mode flag (specified)    & Not yet implemented                                          &                                                                                 \\
		Blanking-only VRAM access (specified)              & Not yet implemented                                          & Not enforced                                                                    \\
		Forced blank, graphics disable (specified)         & Not yet implemented                                          & GD stored only                                                                  \\
		Mode 2                                             & Prototype                                                    & Proposed; Chapter~\ref{ch:mode2}                                                \\
		Mode 2 status layer (background 3)                 & Prototype                                                    & Proposed; Section~\ref{sec:m2-status-layer}                                     \\
		Mode 2 2-bit isometric tiles                       & Prototype                                                    & Proposed; Section~\ref{sec:m2-two-bit}                                          \\
		Mode 3                                             & Prototype                                                    & Proposed; Chapter~\ref{ch:mode3}                                                \\
		Isometric shading, shadow sprites                  & Prototype                                                    & Proposed; Section~\ref{sec:m2-shading}                                          \\
		2-bit orthogonal tiles                             & Mode 3 and Mode 2 only                                       & Not yet implemented in Mode 1                                                   \\
		\bottomrule
	\end{tabular}
	\caption{Implementation Status}
	\label{tab:status}
\end{table}

\section{Discrepancies}
\label{sec:status-discrepancies}

The following are known differences between this manual's source, the simulator, and its other sources. Each is a
defect to be fixed in the simulator or a decision to be made; the manual describes the intended behavior and notes the
difference.

\begin{enumerate}
	\item \textbf{Register addresses.} The wiki puts the sprite tile data base at offset \texttt{0x0012} and the status at
	      \texttt{0x0013}. The simulator puts the three background palette configuration registers at
	      \texttt{0x0012}--\texttt{0x0014}, the sprite tile data base at \texttt{0x0016} and the status at \texttt{0x0017}.
	      This manual follows the simulator.
	\item \textbf{Palette boundary} (fixed). The simulator used to decode palette offsets \texttt{0x6000}--\texttt{0x6100}
	      inclusive, one past the end of its 256 entries; it now stops at \texttt{0x60FF} (Section~\ref{sec:palette-edge}).
	\item \textbf{Mode 1 layer order} (fixed with the sprite prototype). The simulator used to show background 3 in front of
	      background 1; it now uses the intended order (Section~\ref{sec:m1-priority}).
	\item \textbf{Mode 1 palette sum} (fixed with the sprite prototype). The simulator added the palette offset and the
	      tile value without wrapping, so a large offset could stop the simulator in a debug build; it now wraps, as
	      Mode 2 and Section~\ref{sec:palette-select} do.
	\item \textbf{Status byte order} (fixed). The simulator read the Status register with its bytes swapped relative to every
	      other register; it no longer does.
	\item \textbf{Vertical sync position} (fixed). The interrupt used to be asserted at line 0; it is now asserted at line
	      224, the start of vertical blanking (Chapter~\ref{ch:timing}).
	\item \textbf{Mode 1 picture offset.} The simulator's Mode 1 shows the first 8 pixels of every line blank and draws the
	      tilemap 8 pixels to the right of its position, so the last tile column is never shown. The first two groups of a line
	      are flushed before any data has been fetched (Section~\ref{sec:m1-scrolling}).
	\item \textbf{Clock accounting.} The simulator polls every device once per \texttt{CLKI} cycle, which the video device counts
	      as two master clocks (Section~\ref{sec:cpu-clock}); the terminal device counts one poll as one master clock, so its
	      baud rate is off by a factor of two.
	\item \textbf{VRAM access during rendering.} The specification ignores accesses to video memory and the sprite table outside
	      vertical blanking; the simulator allows them (Chapter~\ref{ch:memory-map}).
\end{enumerate}

\section{Open Design Questions}

\subsection{General}

\begin{itemize}
	\item \textbf{Interrupts.} Specified (Section~\ref{sec:interrupts}): vertical sync on the \activelow{NMI} with no
	      acknowledgement, line compare on \activelow{IRQ4}, held until a write to \texttt{LINE}, both enabled from the Interrupt
	      Control register. The simulator still pulses vertical sync on level 4 and line compare on level 3 for one cycle,
	      with no enable for vertical sync. A one-cycle NMI pulse is one \texttt{CLKI} cycle, which the CPU's edge detector accepts.
	\item \textbf{The wiki.} The project wiki's register map (the sprite tile data base and status register offsets) disagrees with
	      this manual and the simulator, and should be updated from this manual.
	\item \textbf{Version and region.} \texttt{VER} is 1 and the region is fixed when the device is made; the simulator should
	      let both be configured at run time.
	\item \textbf{Package and signals.} The package, the pin assignment and the signal timing of Chapters~\ref{ch:package} and~\ref{ch:signal-timing}
	      are a first sketch, to be confirmed: the 52-pin PLCC; a chip select instead of the segment address bits; \texttt{CLKI} as
	      the only clock; \activelow{BACK}, \activelow{NMI} and \activelow{IRQ4} as open-drain outputs; no access-type or protection
	      inputs; a minimum reset of six cycles; the composite sync widths (a horizontal pulse of 100 master clocks; vertical sync
	      in three lines, 245 to 247, one long pulse in each, no equalising pulses); and the converter and supply pins. The video memory,
	      palette and sprite table are inside the package. Electrical characteristics (levels, delays, frequency range, current,
	      thermal) are not yet specified, as in the \hwmanual{}.
	\item \textbf{DMA.} A DMA transfer into video memory is subject to the same rule as the CPU
	      (Section~\ref{sec:memory-vram-access}), and DMA blocks interrupts, the \activelow{NMI} included, so programs keep
	      transfers short and away from line 224 (Section~\ref{sec:interrupts}).
\end{itemize}

\subsection{Mode 2}

\begin{itemize}
	\item \textbf{Whether to approve it.} Mode 2 narrows the device to one kind of picture. That is its purpose, but it is a
	      decision, not a default.
	\item \textbf{Sprites.} With no free fetch slot, Mode 2 sprites come from a line buffer filled in the blanking
	      interval; the prototype shows that 32 sprites per line of either size fit (Chapter~\ref{ch:sprites}). Whether
	      32 is enough for a vehicle-heavy game, and whether to pay for a second line buffer to double it, is open.
	\item \textbf{The status layer.} Prototyped with a queue of four groups (Section~\ref{sec:m2-status-layer}). Open: whether
	      it should scroll, whether its tiles can be flipped, and whether to use its slots for a third isometric background
	      instead.
	\item \textbf{2-bit tiles.} Prototyped (Sections~\ref{sec:m2-two-bit} and~\ref{sec:flat-two-bit}). Left to the program:
	      the palette size with \texttt{T2}. Open: the Mode 1 implementation.
	\item \textbf{Mode 3.} Prototyped (Chapter~\ref{ch:mode3}). Open: whether the orthogonal backgrounds should have a priority bit to
	      draw in front of sprites, and whether another arrangement of orthogonal and isometric backgrounds is wanted.
	\item \textbf{Shading.} Prototyped (Section~\ref{sec:m2-shading}). Open: whether darkening should have more than one level,
	      and whether a shadow pixel should darken sprites behind it as well as backgrounds.
	\item \textbf{The tilemap entry cache.} Measured and left out (Section~\ref{sec:m2-entry-cache}); worth revisiting if
	      the spare slots are wanted for something else.
	\item \textbf{Rotation of the view.} Ruled out as a register; to show another facing a program needs another set of
	      tiles and map.
	\item \textbf{The symmetric diamond.} The footprint of Table~\ref{tab:m2-footprint} is one pixel narrower at the top than the
	      bottom. A conventionally symmetric diamond does not give a bijection (some texels would be shown twice and
	      some never). The manual keeps the footprint that does. The slicing tool accepts any picture, but diamonds that
	      are to tessellate must be drawn to that footprint.
	\item \textbf{The fixed tilemap size.} A 32 by 32 map is 512 by 512 texels; a transport game would want a larger map and
	      would page it by rewriting the tilemap.
\end{itemize}
